\documentclass{article}
\usepackage[T1]{fontenc}
\usepackage{lmodern}
\usepackage{graphicx}
\usepackage{amsmath,amssymb}
\usepackage[a4paper,margin=2cm]{geometry}
\usepackage[absolute,overlay]{textpos}
\setlength{\TPHorizModule}{1mm}
\setlength{\TPVertModule}{1mm}
\usepackage[indonesian]{babel}
\usepackage{hyperref}
\setlength{\emergencystretch}{2em}
% unit: Halaman 36

\begin{document}

% === Halaman 36 (python/native) ===

36

Aplikasi Cipher Alir

• Cipher alir cocok untuk mengenkripsi aliran data yang kontinu melalui 

saluran komunikasi, misalnya:

   1. Mengenkripsi data pada saluran yang menghubungkan  antara dua 

buah komputer.

 2. Mengenkripsi suara pada jaringan telepon mobile GSM.

• Alasan: jika bit cipherteks yang diterima mengandung kesalahan, maka hal 

ini hanya menghasilkan satu bit kesalahan pada waktu dekripsi, karena tiap 

bit plainteks ditentukan hanya oleh satu bit cipherteks.

• Selain itu, alasannya karena cipher alir itu komputasinya cepat (dibutuhkan 

untuk real-time processing)

\end{document}
